\documentclass[
	12pt,				% tamanho da fonte
	openright,			% capítulos começam em pág ímpar (insere página vazia caso preciso)
	oneside,			% para impressão frente. Oposto a twoside
	a4paper,			% tamanho do papel. 
	% -- opções da classe abntex2 --
	chapter=TITLE,		% títulos de capítulos convertidos em letras maiúsculas
	% -- opções do pacote babel --
	english,			% idioma adicional para hifenização
	french,				% idioma adicional para hifenização
	spanish,			% idioma adicional para hifenização
	brazil				% o último idioma é o principal do documento
	]{abntex2}
\usepackage{ufr/ufr} % usando package da ufr.

\usepackage{listings}     
\usepackage{lstautogobble}  % Fix relative indenting
\usepackage{color}          % Code coloring
\usepackage{zi4}            % Nice font

\definecolor{bluekeywords}{rgb}{0.13, 0.13, 1}
\definecolor{greencomments}{rgb}{0, 0.5, 0}
\definecolor{redstrings}{rgb}{0.9, 0, 0}
\definecolor{graynumbers}{rgb}{0.5, 0.5, 0.5}

\usepackage{listings}
\lstset{
    autogobble,
    columns=fullflexible,
    showspaces=false,
    showtabs=false,
    breaklines=true,
    showstringspaces=false,
    breakatwhitespace=true,
    escapeinside={(*@}{@*)},
    commentstyle=\color{greencomments},
    keywordstyle=\color{bluekeywords},
    stringstyle=\color{redstrings},
    numberstyle=\color{graynumbers},
    basicstyle=\ttfamily\footnotesize,
    frame=l,
    framesep=12pt,
    xleftmargin=12pt,
    tabsize=4,
    captionpos=b
}

\usepackage{caption}

% \usepackage[color=red,showframe]{fgruler}

% ---
% compila o indice
% ---
\makeindex
% ---

% ---
% Este é o documento principal do template para Trabalhos Acadêmicos da Universidade Federal de Rondonópolis. Esse template usa o AbnTeX2 , que possui diversos comandos para elaboração de um documento acadêmico.
% Você pode modificar esse documento se achar necessário, mas recomendamos que evite fazer isso.
% ---
% ----
% Início do documento
% ----
\begin{document}

% Seleciona o idioma do documento (conforme pacotes do babel)
%\selectlanguage{english}
\selectlanguage{brazil}

% Retira espaço extra obsoleto entre as frases.
\frenchspacing 
% ----------------------------------------------------------
% ELEMENTOS PRÉ-TEXTUAIS
% ----------------------------------------------------------
\pretextual
\imprimircapa
\imprimirfolhaderosto
%\imprimirfolhaderosto* % se for incluir ficha catalográfica

% Isto é um exemplo de Ficha Catalográfica, ou ``Dados internacionais de
% catalogação-na-publicação''. Você pode utilizar este modelo como referência. 
% Porém, provavelmente a biblioteca da sua universidade lhe fornecerá um PDF
% com a ficha catalográfica definitiva após a defesa do trabalho. Quando estiver
% com o documento, salve-o como PDF no diretório do seu projeto e substitua todo
% o conteúdo de implementação deste arquivo pelo comando abaixo:
%


% \begin{fichacatalografica}
%     \includepdf{fig_ficha_catalografica.pdf}
% \end{fichacatalografica}



% \begin{fichacatalografica}
% 	\sffamily
% 	\vspace*{\fill}					% Posição vertical
% 	\begin{center}					% Minipage Centralizado
% 	\fbox{\begin{minipage}[c][8cm]{13.5cm}		% Largura
% 	\small
% 	\imprimirautor
% 	%Sobrenome, Nome do autor
	
% 	\hspace{0.5cm} \imprimirtitulo  / \imprimirautor. --
% 	\imprimirlocal, \imprimirdata-
	
% 	\hspace{0.5cm} \thelastpage p. : il. (algumas color.) ; 30 cm.\\
	
% 	\hspace{0.5cm} \imprimirorientadorRotulo~\imprimirorientador\\
	
% 	\hspace{0.5cm}
% 	\parbox[t]{\textwidth}{\imprimirtipotrabalho~--~\imprimirinstituicao,
% 	\imprimirdata.}\\
	
% 	\hspace{0.5cm}
% 		1. Palavra-chave1.
% 		2. Palavra-chave2.
% 		2. Palavra-chave3.
% 		I. Orientador.
% 		II. Universidade xxx.
% 		III. Faculdade de xxx.
% 		IV. Título 			
% 	\end{minipage}}
% 	\end{center}
% \end{fichacatalografica}
% \begin{errata}
% Elemento opcional da {NBR14724:2011}. Exemplo:

% \vspace{\onelineskip}

% FERRIGNO, C. R. A. \textbf{Tratamento de neoplasias ósseas apendiculares com
% reimplantação de enxerto ósseo autólogo autoclavado associado ao plasma
% rico em plaquetas}: estudo crítico na cirurgia de preservação de membro em
% cães. 2011. 128 f. Tese (Livre-Docência) - Faculdade de Medicina Veterinária e
% Zootecnia, Universidade de São Paulo, São Paulo, 2011.

% \begin{table}[htb]
% \center
% \footnotesize
% \begin{tabular}{|p{1.4cm}|p{1cm}|p{3cm}|p{3cm}|}
%   \hline
%    \textbf{Folha} & \textbf{Linha}  & \textbf{Onde se lê}  & \textbf{Leia-se}  \\
%     \hline
%     1 & 10 & auto-conclavo & autoconclavo\\
%    \hline
% \end{tabular}
% \end{table}

% \end{errata}
% Isto é um exemplo de Folha de aprovação, elemento obrigatório da NBR
% 14724/2011 (seção 4.2.1.3). Você pode utilizar este modelo até a aprovação do trabalho. Após isso, substitua todo o conteúdo deste arquivo por uma imagem da página assinada pela banca com o comando abaixo:
%
% \begin{folhadeaprovacao}
% \includepdf{folhadeaprovacao_final.pdf}
% \end{folhadeaprovacao}

\begin{folhadeaprovacao}% 
    \begin{center}
    
        {\bfseries\imprimirautor}
        
        \vspace{3\baselineskip}
        
        {\bfseries\imprimirtitulo}
        
        \vspace{3\baselineskip}
    
    \end{center}

    %preambulo
    \hspace{.4\textwidth}
    \begin{minipage}{.5\textwidth}
        \imprimirpreambulo
    \end{minipage}%
    
   
    \vspace{3\baselineskip}

    \begin{center}
        \imprimirdataaprovacao
    \end{center}  
    
    \vspace{3\baselineskip}
    \begin{center}
        \textbf{Banca Examinadora:}
    \end{center}
   \vspace*{-1cm}
   
   \assinatura{\textbf{\imprimirorientador} \\Orientador \\ UFR} 
   \assinatura{\textbf{Nome do Co-orientador} \\ Co-orientador \\ UFR ou outra instituição }
   
   \assinatura{\textbf{Nome Externo} \\ Examinador(a) \\ UFR ou outra instituição}
   % \assinatura{\textbf{Nome Externo} \\ Examinador(a) \\ UFR ou outra instituição}
   

\end{folhadeaprovacao}
\begin{dedicatoria}
   Dedicatória. Elemento opcional. Frase por meio da qual o autor presta homenagem ou dedica o trabalho a alguém. O título Dedicatória não deve aparecer no início da folha. Se não desejar escrever uma dedicatória, delete este elemento.
\end{dedicatoria}

% ---
% Agradecimentos
% ---
\begin{agradecimentos}

\textit{Elemento opcional. Menção que o autor faz a pessoas e/ou instituições que colaboraram de maneira relevante na elaboração do trabalho. Se não desejar escrever agradecimentos, delete esta página. }

Ao professor fulano, pela...

Ao Senhor cicrano, pela...

A todos que direta ou indiretamente colaboraram na execução deste trabalho.
\end{agradecimentos}
% ---
% ---
% Epígrafe
% ---
\begin{epigrafe}
		Não vos amoldeis às estruturas deste mundo,  mas transformai-vos pela renovação da mente a fim de distinguir qual é a vontade de Deus: o que é bom, o que Lhe é agradável, o que é perfeito.

        (Bíblia Sagrada, Romanos 12, 2)	
\end{epigrafe}
% ---

% ---
% RESUMOS
% ---
% resumo em português

\begin{resumo}
 De acordo com a ABNT NBR 6028:2021, o resumo informativo deve ressaltar o objetivo, o método, os resultados e as conclusões do documento. Ele deve ser composto de uma sequência de frases concisas e afirmativas. Convém usar o verbo na terceira pessoa do singular. O texto do resumo deve ser digitado em um parágrafo único, justificado. O espaçamento entre linhas é simples e o tamanho da fonte é 12. Deve conter de 150 a 500 palavras. As palavras-chave devem figurar logo abaixo do resumo, antecedidas da expressão Palavras-chave, seguida de dois-pontos, separadas entre si por ponto e vírgula e finalizadas por ponto. Devem ser grafadas com as iniciais em letra minúscula, com exceção dos substantivos próprios e nomes científicos.
 
\textbf{Palavras-chave}: latex; abntex; editoração de texto.

\end{resumo}


% resumo em inglês
\begin{resumo}[Abstract]
 \begin{otherlanguage*}{english}
   \lipsum[150]

   \vspace{\onelineskip}
 
   \noindent 
   \textbf{Keywords}: latex; abntex; text editoration.
 \end{otherlanguage*}
\end{resumo}

% Se quiser, é possível adicionar outros idiomas para resumo, bastando utilizar os exemplos abaixo.  

% % resumo em francês 
% \begin{resumo}[Résumé]
% \SingleSpacing
%  \begin{otherlanguage*}{french}
%     Il s'agit d'un résumé en français.
 
%    \textbf{Mots-clés}: latex. abntex. publication de textes.
%  \end{otherlanguage*}
% \end{resumo}

% % resumo em espanhol
% \SingleSpacing
% \begin{resumo}[Resumen]
%  \begin{otherlanguage*}{spanish}
%    Este es el resumen en español.
  
%    \textbf{Palabras clave}: latex. abntex. publicación de textos.
%  \end{otherlanguage*}
% \end{resumo}
% ---
% Aqui vão as listas diversas, na ordem que se encontram. Caso não haja alguma das listas você pode simplesmente eliminá-la removendo o código relacionado à ela.

% ---
% inserir lista de ilustrações
% ---
\pdfbookmark[0]{\listfigurename}{lof}
\listoffigures*
\cleardoublepage
% ---

% ---
% inserir lista de quadros
% ---
\pdfbookmark[0]{\listofquadrosname}{loq}
\listofquadros*
\cleardoublepage
% ---

% ---
% Inserir lista de tabelas
% ---
\pdfbookmark[0]{\listtablename}{lot}
\listoftables*
\cleardoublepage
% ---

% ---
%%% Inserir lista de abreviaturas e siglas %%% 
% Para inserir uma nova abreviatura/sigla, basta usar o formato abaixo
% ---
\begin{siglas}
  \item[ABNT] Associação Brasileira de Normas Técnicas
  \item[abnTeX] ABsurdas Normas para TeX
\end{siglas}
% ---

% ---
%%% Inserir lista de símbolos %%%
% Para inserir uma nova abreviatura/sigla, basta usar o formato abaixo
% ---
\begin{simbolos}
  \item[$ \Gamma $] Letra grega Gama
  \item[$ \Lambda $] Lambda
  \item[$ \zeta $] Letra grega minúscula zeta
  \item[$ \in $] Pertence
\end{simbolos}
% ---

\imprimirsumario

\textual %NÃO REMOVER 
\pagestyle{simple} %NÃO REMOVER 


% ----------------------------------------------------------
% Introdução (exemplo de capítulo sem numeração, mas presente no Sumário)
% ----------------------------------------------------------
\chapter{Introdução}

A introdução é o primeiro elemento textual e contém alguns itens importantes do projeto de pesquisa: tema, questões de pesquisa, objetivos e justificativa (sucinta). Deve situar o autor da pesquisa em relação ao que irá estudar, apresentando em linhas gerais como chegou ao tema e como pretende desenvolvê-lo em sua pesquisa. Ela deve se encerrar apresentando ao leitor a organização retórica de seu trabalho, ou seja, as partes que compõem o TCC.

\section{Olá mundo}

Segundo Machado, Lousada e Abreu-Tardelli (2005, p. 83), “a introdução pode ser vista como um trailer do que o leitor verá no seu trabalho, nem mais nem menos”. É uma seção que deve levar o leitor a querer ler o trabalho, seduzindo-o.

\subsection{Olá mundo}

Uma dica útil dada pelas autoras é apresentar inicialmente o “objeto” sobre o qual trata a pesquisa em um relato de como você chegou ao tema, quais os motivos mais relevantes, as buscas que efetuou, as decisões tomadas e as teorias que foi selecionando ao longo dessa busca. 

\subsubsection{Olá mundo}

Aqui serão dadas indicações gerais para a apresentação gráfica de seu trabalho, contudo, você pode consultar a NBR 14724:2011 para obter mais informações sobre a apresentação de trabalhos acadêmicos.

\subsubsubsection{Olá mundo}

Todo o texto deve ser digitado em espaço 1,5 cm, exceto: citações de mais de três linhas, notas de rodapé, referências, legendas e fontes das ilustrações e das tabelas que devem ser digitados em espaço simples. As referências, ao final do trabalho, devem ser separadas entre si por um espaço simples em branco.

\begin{citacao}
    Esta é uma citação, \lipsum[20]
\end{citacao}

As margens da página devem ser de 3 cm nas margens esquerda e superior e 2 cm nas margens direita e inferior.

Os títulos dos capítulos (seções primárias, secundárias, etc.) devem ser digitados após a sua numeração (indicação de seção), separados por um espaço. O texto deve iniciar em outra linha, separado por um espaço entrelinhas de 1,5. 

As seções primárias iniciam-se em nova página e são grafadas em caixa alta e negrito. As seções secundárias são grafadas em negrito com apenas a primeira letra maiúscula. As seções terciárias não são grafadas em negrito. Escreva um título criativo consoante o arcabouço teórico e seu plano de trabalho constante no sumário. Utilize junto a este template a NBR 6024:2012 – Numeração progressiva das seções de um documento.

A fonte utilizada no texto é Arial ou Times, tamanho 12, excetuando-se citações com mais de três linhas, notas de rodapé, paginação, legendas e fontes das ilustrações e das tabelas, que devem ser em tamanho 11. 

As citações diretas com mais de três linhas devem ser destacadas com recuo de 4 cm da margem esquerda, fonte tamanho 11 e sem as aspas. Consulte sempre a norma específica para citações, NBR 10520:2002.

As páginas pré-textuais (todas que precedem a Introdução) devem ser contadas, mas não numeradas, exceto a capa e página da Ficha Catalográfica. A numeração deve figurar a partir da Introdução, em algarismos arábicos, no canto superior direito da folha, a 2 cm da borda superior, ficando o último algarismo a 2 cm da borda direita da folha, fonte 11.


% \input{tabelas/tabela1.tex}
%\chapter{Trabalhando com Figuras e códigos}


\section{Figuras}

A ABNT demanda que, ao adicionar figuras, devamos adicioná-las o mais próximo possível da sua primeira citação. O template (ainda) não controla isso. Para adicionar figuras, fazemos como abaixo \ref{fig:figura-logo}

\begin{figure}
    \centering
    \includegraphics{ufr/images/logo.jpeg} % localização da figura
    \caption{Legenda}
    \label{fig:figura-logo}
\end{figure}
% ---


%% ---
% primeiro capitulo de Resultados
% ---
\chapter{Lectus lobortis condimentum}
% ---

% ---
\section{Vestibulum ante ipsum primis in faucibus orci luctus et ultrices
posuere cubilia Curae}
% ---

\cite{}

\lipsum[21-22]

%\chapter{Revisão de Literatura}\label{cap_trabalho_academico}

\section{Quadros}
Este modelo vem com o ambiente \texttt{quadro} e impressão de Lista de quadros 
configurados por padrão. Verifique um exemplo de utilização:

\begin{quadro}[htb]
\caption{\label{quadro_exemplo}Exemplo de quadro}
\begin{tabular}{|c|c|c|c|}
	\hline
	\textbf{Pessoa} & \textbf{Idade} & \textbf{Peso} & \textbf{Altura} \\ \hline
	Marcos & 26    & 68   & 178    \\ \hline
	Ivone  & 22    & 57   & 162    \\ \hline
	...    & ...   & ...  & ...    \\ \hline
	Sueli  & 40    & 65   & 153    \\ \hline
\end{tabular}
\fonte{Autor.}
\end{quadro}

\textbf{Dica:} Quadros apresentam dados não numéricos. Tabelas apresentam dados numéricos.

\subsection{Quadro 2}

Este parágrafo apresenta como referenciar o quadro no texto, requisito
obrigatório da ABNT. 
Primeira opção, utilizando \texttt{autoref}: Ver o \autoref{quadro_exemplo}. 
Segunda opção, utilizando  \texttt{ref}: Ver o Quadro \ref{quadro_exemplo}.

\section{Tabelas}


%% ---
% Conclusão
% ---
\chapter{Conclusão}
% ---

\lipsum[31-33]

\chapter{Escrevendo referências bibliográficas}

sjsjsjsjsjsj \cite{ersi-rj}

\postextual

\bibliography{bibliografia}

\begin{apendicesenv}
        % % ----------------------------------------------------------
\chapter{Quisque libero justo}
% ----------------------------------------------------------

\lipsum[50]
         
        % % ----------------------------------------------------------
\chapter{Nullam elementum urna vel imperdiet}
% ----------------------------------------------------------
\lipsum[55-57]
    \end{apendicesenv}
    % ---

\begin{anexosenv}
    % % ---
\chapter{Morbi ultrices rutrum lorem.}
% ---
\lipsum[30]

    % % ---
\chapter{Cras non urna sed feugiat cum sociis natoque penatibus et 
parturient montes nascetur ridiculus mus}
% ---

\lipsum[31]

\end{anexosenv}

\end{document}